\section*{Chapter \ref{ch:b3:lithosphere-thermal-models} --- Thermal Models of the Lithosphere}

\begin{solution}{exo:b3:lithosphere-thermal-models:1}
$L^2/\kappa = \qty{e12}{s}$, \qty{e14}{s} and \qty{e16}{s}: 0.032, 3.2 and
\qty{317}{Myr}. A kilometre forgets its temperature in tens of thousands of
years, a hundred kilometres in hundreds of millions.
\end{solution}

\begin{solution}{exo:b3:lithosphere-thermal-models:2}
$u = \qty{e-3}{m}/\qty{3.156e7}{s} = \qty{3.17e-11}{m.s^{-1}}$, so
$\mathrm{Pe} = 3.17\times10^{-11}\times\num{3e4}/10^{-6} = 0.95$. Advection and
conduction are of the same importance: the rising rock carries heat upwards
and steepens the geotherm near the surface, and conduction alone would
underestimate the gradient.
\end{solution}

\begin{solution}{exo:b3:lithosphere-thermal-models:3}
$\sqrt{\kappa t} = \sqrt{\num{0.8e-6}\times5\times\num{3.156e13}} =
\qty{11236}{m}$, $\eta = 10000/(2\times11236) = 0.445$, and $T = 1450\times0.471
= \qty{683}{\celsius}$.
\end{solution}

\begin{solution}{exo:b3:lithosphere-thermal-models:4}
$490/\sqrt1 = 490$, $490/5 = 98$ and $490/10 = \qty{49}{mW.m^{-2}}$. The probe
would most likely underestimate the first: young crust is cracked and open to
circulating sea water, which carries part of the heat away by hot springs.
\end{solution}

\begin{solution}{exo:b3:lithosphere-thermal-models:5}
$z = 2\times1.163\sqrt{\kappa t}$: with $\sqrt{\kappa t} = 15.89$, 31.78 and
\qty{50.25}{km}, $z = 37.0$, 73.9 and \qty{116.9}{km}. The cooled layer doubles
from 10 to \qty{40}{Myr}; by \qty{100}{Myr} it is already deeper than a
\qty{95}{km} plate.
\end{solution}

\begin{solution}{exo:b3:lithosphere-thermal-models:6}
$C$ is proportional to $\alpha$: $\alpha = \num{3.1e-5}\times350/370.5 =
\qty{2.93e-5}{K^{-1}}$. $C$ is also proportional to $T_m - T_0$ (a mantle at
$1450\times350/370.5 = \qty{1370}{\celsius}$) and to $\sqrt\kappa$ (a
diffusivity $0.8\times(350/370.5)^2 = \qty{0.71e-6}{m^2.s^{-1}}$): depth alone
cannot separate them, which is why heat flow is used as well.
\end{solution}

\begin{solution}{exo:b3:lithosphere-thermal-models:7}
$3.1\times1450/\num{95e3} = \qty{47.3}{mW.m^{-2}}$ and $3.1\times1350/\num{125e3}
= \qty{33.5}{mW.m^{-2}}$. The half-space reaches them at $(490/47.3)^2 =
\qty{107}{Myr}$ and $(490/33.5)^2 = \qty{214}{Myr}$: the thin, hot plate holds
the heat flow of old floor at \qty{47}{mW.m^{-2}} from about \qty{100}{Myr}
on, whereas the thicker, colder one would let it fall below
\qty{40}{mW.m^{-2}}, the half-space value at \qty{150}{Myr}.
\end{solution}

\begin{solution}{exo:b3:lithosphere-thermal-models:8}
$A_0h_r = 0.4\times60 = \qty{24}{mW.m^{-2}}$, so $A_0 =
\num{24e-3}/\num{e4} = \qty{2.4}{\micro W.m^{-3}}$ and $q_m = 60 - 24 =
\qty{36}{mW.m^{-2}}$. At \qty{35}{km}, $T = \num{36e-3}\times\num{35e3}/2.5 +
\num{2.4e-6}\times\num{e8}/2.5\times(1 - e^{-3.5}) = 504 + 96\times0.970 =
\qty{597}{\celsius}$.
\end{solution}

\begin{solution}{exo:b3:lithosphere-thermal-models:9}
$h_r = \Delta q_0/\Delta A_0 = \num{18e-3}/\num{1.8e-6} = \qty{10}{km}$, and
$q_m = q_0 - h_rA_0 = 66 - 30 = \qty{36}{mW.m^{-2}}$, since $h_rA_0 =
\num{e4}\times\qty{3.0e-6}{W.m^{-3}} = \qty{30}{mW.m^{-2}}$.
\end{solution}

\begin{solution}{exo:b3:lithosphere-thermal-models:10}
(a) $Q(t) = \int_0^tk(T_m - T_0)/\sqrt{\pi\kappa t'}\,\mathrm dt' =
2k(T_m - T_0)\sqrt{t/(\pi\kappa)}$. (b) The heat lost is the heat content
removed from the column, $Q = \rho_mc\int_0^\infty(T_m - T)\,\mathrm dz$, and
isostasy gives $(\rho_m - \rho_w)(d - d_r) = \rho_m\alpha\int_0^\infty(T_m -
T)\,\mathrm dz = \alpha Q/c$. (c) $c = 3.1/(3300\times\num{0.8e-6}) =
\qty{1174}{J.kg^{-1}.K^{-1}}$; at \qty{40}{Myr}, $t = \qty{1.262e15}{s}$ and
$Q = 2\times3.1\times1450\times\sqrt{\num{1.262e15}/(\pi\times\num{0.8e-6})} =
\qty{2.015e14}{J.m^{-2}}$; $d - d_r = \num{3.1e-5}\times\num{2.015e14}/
(1174\times2270) = \qty{2343}{m}$, which is $370.5\sqrt{40}$.
\end{solution}

\begin{solution}{exo:b3:lithosphere-thermal-models:11}
The $n = 3$ term is $3104\times(8/\pi^2)\times\tfrac19e^{-9t/36.2} =
280\,e^{-9t/36.2}$ metres: \qty{218}{m} at \qty{1}{Myr}, \qty{23}{m} at
\qty{10}{Myr} and \qty{1.9}{m} at \qty{20}{Myr}. It falls to \qty{10}{m} when
$e^{-9t/36.2} = 10/280$, at $t = (36.2/9)\ln28 = \qty{13.4}{Myr}$.
\end{solution}

\begin{solution}{exo:b3:lithosphere-thermal-models:12}
(a) For $z \gg h_r$, $e^{-z/h_r} \to 0$. (b) $q_m = \qty{24}{mW.m^{-2}}$,
$A_0h_r^2/k = \num{16e-3}\times\num{e4}/2.5 = \qty{64}{K}$, and $z = (1300 -
64)\times2.5/\num{24e-3} = \qty{128.8}{km}$. (c) $q_m = 2.5\times1236/\num{250e3}
= \qty{12.4}{mW.m^{-2}}$: the \omterm{def:b3:lithosphere-thermal-models:reduced}{reduced heat flow} must be about halved, as under a
craton whose thick, cold root lets little heat in from below.
\end{solution}

\begin{solution}{pb:b3:lithosphere-thermal-models:1}
\textbf{1.} $\sqrt{\num{0.8e-6}\times\num{3.156e14}} = \qty{15889}{m}$ and
$\sqrt{\num{0.8e-6}\times\num{1.2624e15}} = \qty{31778}{m}$.

\textbf{2.} $\eta = 20000/31778 = 0.6294$ and $20000/63556 = 0.3147$:
$T = 1450\times0.6266 = \qty{909}{\celsius}$ at \qty{10}{Myr} and
$1450\times0.3437 = \qty{498}{\celsius}$ at \qty{40}{Myr}.

\textbf{3.} $z = 2\times0.7174\sqrt{\kappa t}$: \qty{22.8}{km} and
\qty{45.6}{km}. It reaches \qty{60}{km} when $\sqrt{\kappa t} = 60/1.4348 =
\qty{41.82}{km}$, at $t = (\num{41.82e3})^2/\num{0.8e-6} = \qty{2.186e15}{s}$,
\qty{69.3}{Myr}.

\textbf{4.} $q = 3.1\times1450/(\sqrt\pi\times15889) = \qty{159.6}{mW.m^{-2}}$ at
\qty{10}{Myr} and \qty{79.8}{mW.m^{-2}} at \qty{40}{Myr}; $C_q = 159.6\sqrt{10}
= \qty{505}{mW.m^{-2}.Myr^{1/2}}$.

\textbf{5.} $u = \qty{7.92e-10}{m.s^{-1}}$, $\mathrm{Pe} = 7.92\times10^{-10}
\times\num{95e3}/\num{0.8e-6} = 94$. Heat travels sideways with the plate a
hundred times faster than it diffuses sideways, so each column cools by
vertical conduction only.

\textbf{6.} $(\num{95e3})^2/\num{0.8e-6} = \qty{1.13e16}{s}$, \qty{357}{Myr}.
At \qty{150}{Myr}, $\sqrt{\kappa t} = \qty{61.5}{km}$ and the cooled layer, a
few times that, reaches far below \qty{95}{km}: a half-space model of so old a
floor assumes that cold has reached depths where the mantle convects.

\textbf{7.} $(\rho_m - \rho_w)(d - d_r) = \rho_m\alpha\int(T_m - T)\,\mathrm dz$,
the masses of the ridge column and of the old column being equal above a
compensation depth. $T_m - T = T_m\operatorname{erfc}\eta$ falls below
\qty{0.5}{\percent} of $T_m$ for $\eta > 2$, so the deep part adds nothing.

\textbf{8.} $C = 2\times3300\times\num{3.1e-5}\times1450/2270\times
\sqrt{\num{0.8e-6}\times\num{3.156e13}/\pi} = 0.13069\times2834.8 =
\qty{370.5}{m.Myr^{-1/2}}$; $d = 2500 + 370.5\sqrt{10} = \qty{3672}{m}$ and
$2500 + 370.5\sqrt{40} = \qty{4843}{m}$.

\textbf{9.} The model gives $2500 + 370.5\sqrt{42.5} = \qty{4915}{m}$,
\qty{440}{m} deeper than the median. The median includes the load of the
sediment, the seamounts and volcanic plateaus, the swells of hotspots, and
cells on the flanks of fracture zones and margins, all of which make the floor
shallower.

\textbf{10.} $C$ is proportional to $T_m$: $370.5\times1350/1450 =
\qty{344.9}{m.Myr^{-1/2}}$, close to the fitted 350.

\textbf{11.} Sea water circulating through the cracked young crust carries
part of the heat out through hot springs, between the probes.

\textbf{12.} $\tau = (\num{95e3})^2/(\pi^2\times\num{0.8e-6}) =
\qty{1.143e15}{s} = \qty{36.2}{Myr}$; the amplitude is $3300\times
\num{3.1e-5}\times1450\times\num{95e3}/4540 = \qty{3104}{m}$.

\textbf{13.} $e^{-40/36.2} = 0.33124$, $e^{-360/36.2}/9 = \num{5.3e-6}$ and the
$n = 5$ term is smaller still; the bracket is $1 - 0.81057\times0.33125 =
0.73150$, and $d = 2500 + 3104\times0.73150 = \qty{4771}{m}$.

\textbf{14.} By the $n = 3$ term, $3104\times0.81057\times\num{5.3e-6} =
\qty{0.013}{m}$: at \qty{40}{Myr}, $e^{-9t/\tau}$ is $10^{-4}$ of
$e^{-t/\tau}$, and only the first mode is left.

\textbf{15.} At \qty{40}{Myr}, $4843 - 4771 = \qty{73}{m}$; at
\qty{50}{Myr}, $2500 + 370.5\sqrt{50} = 5120$ against $2500 + 3104\times(1 -
0.81057\times0.25127) = 4972$: \qty{148}{m}. Interpolating, $40 + 10\times
(100 - 73)/(148 - 73) = \qty{43.6}{Myr}$, about \qty{44}{Myr} (the full series
gives 44.0).

\textbf{16.} $q = (3.1\times1450/\num{95e3})[1 + 2(0.33124 + 0.01204 +
0.00005)] = 47.32\times1.68666 = \qty{79.8}{mW.m^{-2}}$, the same as the
half-space: at \qty{40}{Myr} the base has not yet made itself felt at the
surface.

\textbf{17.} $2500 + 3104 = \qty{5604}{m}$ and $kT_m/L = \qty{47.3}{mW.m^{-2}}$.

\textbf{18.} $d(120) = 2500 + aL\bigl[1 - (8/\pi^2)e^{-120/\tau}\bigr]$, with
$a = \rho_m\alpha T_m/\bigl(2(\rho_m - \rho_w)\bigr) = \qty{32.673}{m/km}$ and
$\tau = L^2/(\pi^2\kappa)$.

\textbf{19.} $L = 3200/\bigl(32.673\,[1 - 0.81057e^{-120/\tau}]\bigr)$. From
$L = \qty{100}{km}$: $\tau = \qty{40.13}{Myr}$, bracket 0.95925, $L =
\qty{102.10}{km}$; then $\tau = 41.83$, $L = 102.67$; $\tau = 42.30$, $L =
102.82$; $\tau = 42.43$, $L = 102.87$; it settles at $L = \qty{102.9}{km}$.

\textbf{20.} $\tau = \qty{42.48}{Myr}$, $aL = \qty{3362}{m}$: $d(40) = 2500 +
3362\times(1 - 0.81057e^{-40/42.48}) = \qty{4799}{m}$, against \qty{4843}{m}
for the half-space and \qty{4800}{m} measured: the plate fits both ages.

\textbf{21.} $\tau = \qty{42.5}{Myr}$; asymptotic depth $2500 + 3362 =
\qty{5862}{m}$.

\textbf{22.} $3362\times0.81057\,e^{-t/42.48} = 100$ gives $t = 42.48\ln27.25 =
\qty{140}{Myr}$.

\textbf{23.} The plate that fits the old floor of the invented ocean is
\textbf{\qty[text-series-to-math=true]{102.9}{km}} thick.
\end{solution}
